\documentclass[10pt,a4paper]{amsart}
\usepackage[margin=19mm]{geometry}
\usepackage{amsmath,amssymb,mathtools}
\usepackage[scr=rsfs,cal=euler]{mathalfa}
\usepackage{xfrac}
\usepackage[unicode,hidelinks,bookmarks=false]{hyperref}
\newcommand{\ie}{i.e.\ }
\newcommand{\cf}{cf.\ }
\newcommand{\resp}{respectively\ }
\newcommand{\Subsection}{\subsection}
\providecommand{\nobreakdash}{\nobreak\hbox{-}}
\setlength{\emergencystretch}{2em}
\multlinegap0pt
\usepackage{amssymb}
\usepackage{mathtools}
\usepackage{csquotes}
\usepackage{xcolor}
\newtheorem{prop}{Proposition}
\newtheorem{rem}{Remark}
\newcommand{\ud}{\mathrm{d}}
\title{\textbf{Expansive solutions and the boundary at infinity for the homogeneous $N$-body problem}}
\author{Diego Berti, Davide Polimeni, Susanna Terracini}
\date{Source: arXiv:2604.14948v1 (CC BY 4.0)}
\begin{document}
\maketitle

\noindent\emph{Source and licence.} \textbf{Expansive solutions and the boundary at infinity for the homogeneous $N$-body problem}. Version arXiv:2604.14948v1 (CC BY 4.0), distributed by arXiv under the Creative Commons Attribution 4.0 International licence, http://creativecommons.org/licenses/by/4.0/, as recorded in the arXiv licence metadata for this exact version and shown on the abstract page https://arxiv.org/abs/2604.14948v1. Full licence notice: \texttt{licenses/arxiv-2604.14948-CC-BY-4.0.txt}.\par\smallskip

\noindent\textit{Editorial scope.} Counted unit: the article's prop prop:asymptotic\_par (source0.tex lines 1776--1785). The article's complete original proof of it is delivered with it (source0.tex lines 1892--2081, one \texttt{proof} environment of 190 lines, which the article places away from the statement). No mathematics is written, repaired or re-derived: the statement and the proof are the article's own lines, in the article's own notation, and no retained line is substituted or re-flowed.  Retained uncounted as the source-local context the proof needs: lines 1810--1813; lines 1825--1840; lines 1843--1851; lines 1869--1872; lines 1882--1885.  Nothing was omitted from any retained span, so the package carries no omission record.  No retained line contains a citation, so the wrapper carries no bibliography and no bibliography entry.  No retained line contains a figure environment, an image-inclusion command or drawing code, so no image is delivered and the package declares no asset dependency.  Editorial formatting only, in addition: an AMS \texttt{amsart} preamble that restates the article's own declarations and macros used by the retained lines, namely the environments \texttt{prop}, \texttt{rem} on separate counters, together with the standard packages those lines need.  The retained environments print in this document as Proposition 1, Remark 1 in order of appearance, whereas the article prints the same items with the numbering of its own full document.  The complete native source of the article as distributed in the arXiv e-print for this exact version is bundled unchanged as \texttt{sources/198568/source0.tex}.
\begin{equation}
\label{eq:mu_f}
\ddot{y}(t) + \frac{\mu}{t^2}y(t) = f(t), \ t> T,
\end{equation}

\begin{rem}
\label{rem:f_gamma}
    Suppose that, for some $q>0$, $f=O(t^{-q})$, for $t\rightarrow+\infty$. Then,
    \[
    \dot y_p(t)=\frac{1}{\theta_+-\theta_-}\left(\theta_+\,t^{\theta_+-1}\int^t\,\sigma^{\theta_-}\,f(\sigma)\,\ud\sigma - \theta_- \, t^{\theta_--1}\int^t \sigma^{\theta_+}\,f(\sigma)\,\ud\sigma\right) =O(t\,|f(t)|)=O(t^{1-q}).
    \]
    Hence, $y_p \in \mathcal{D}_0^{1,2}(T,+\infty)$ if $q>3/2$. 
    
    Moreover, we have  
    \[
    \|y_p(t)\|= O(t^{{\theta_+}+{\theta_-}+1-q})=O(t^{2-q}),
    \]
    for $t\rightarrow+\infty$.

    Since $t^{\theta_+} \notin \mathcal{D}^{1,2}(T,+\infty)$, if $q > 2-\theta_-$, we deduce that every solution $y$ of \eqref{eq:mu_f} lying in $\mathcal{D}^{1,2}(T,+\infty)$ must satisfy $y(t)=O(t^{\theta_-})$, as $t\to+\ \infty$. In the case $\theta_-=\alpha/(2+\alpha)$, then $q > 2-\theta_-$ reads as $q > (4+\alpha)/(2+\alpha)$.
\end{rem}

\begin{prop}[Maximum principle]\label{prop:max_principle}
Let $y\in\mathcal{D}^{1,2}(T,+\infty)$ be such that $
\ddot{y}(t) + \frac{\mu}{t^2}y(t) \ge f(t)$, for  $t>T$, with $\mu <1/4$ and $f \in L^2(T,+\infty;t^2\,\ud t)$.

Then,
\[
y(t) \le y(T)\bigg(\frac{t}{T}\bigg)^{\theta_-} + y_p(t), \ \mbox{ for } \ t\ge T.
\]
\end{prop}

\begin{rem}
\label{rem:corollary}
    A consequence of Proposition \ref{prop:max_principle} and Remark \ref{rem:f_gamma} is that, if $f=O(t^{-q})$, for some $q \ge 2-\theta_-$, and $y\in \mathcal{D}^{1,2}(T,+\infty)$ satisfies $\ddot y(t)+\frac{\mu}{t^2}y(t)\ge f(t)$, for $t >T$, then $y=O(t^{\theta_-})$, for $t\rightarrow+\infty$.
\end{rem}

\begin{equation}
\label{eq:hessian}{
\nabla^2 U(x) = - \alpha\frac{\tilde U(\hat x)}{\|x\|_{\mathcal M}^{2+\alpha}}\mathcal M+\alpha(\alpha+2)\frac{\tilde U(\hat x)}{\|x\|_{\mathcal M}^{2+\alpha}}\mathcal M\hat x\otimes \mathcal M\hat x - (1+2\alpha)\frac{\nabla\tilde U(\hat x)\otimes \mathcal M\hat x}{\|x\|_{\mathcal M}^{2+\alpha}}  +\frac{\nabla^2\tilde  U(\hat x)}{\|x\|_{\mathcal M}^{2+\alpha}},}
\end{equation}

\begin{prop}
\label{prop:asymptotic_par}
    For $\alpha \in (0,2)$, $x \in \mathcal X$, and a normalized minimal central configuration $b_m$, let $\gamma$ be a corresponding parabolic solution having the form $\gamma(t) = \beta b_mt^{\frac{2}{2+\alpha}}+\varphi(t)+x-\beta b_m$, with $\varphi\in\mathcal{D}_0^{1,2}(1,+\infty)$ minimizer of $\mathcal A$, $\beta = \big({\frac{(2+\alpha)^2}{2}U(b_m)}\big)^{\frac{1}{2+\alpha}}$, and $\gamma(1)=x$. 
    
    Then, it holds
    \begin{equation}
    \label{eq:asymp}
    \gamma(t)-\beta b_m\,t^\frac{2}{2+\alpha}\,  =O(t^\frac{\alpha}{2+\alpha}), \ \mbox{ for } \ t \rightarrow+\infty.
    \end{equation}
\end{prop}

\begin{proof}[Proof of Proposition \ref{prop:asymptotic_par}]
    Denoting $r_0(t) = \beta b_m t^{\frac{2}{2+\alpha}}$ and 
    \[
    \psi(t):=\varphi(t)+x-r_0(1),
    \]
    we want to prove that $\|\psi(t)\|_\mathcal{M}=O(t^\frac{\alpha}{2+\alpha})$, as $t \rightarrow+\infty$. We structure the proof into several steps.

    \smallskip
    \emph{Step 1. Intermediate estimates.} \ 
    Define $\eta(t):=\|\psi(t)\|_\mathcal{M}=\left(\sum_i m_i \|\psi_i\|^2\right)^{1/2}$. We have:
\[
\begin{split}
\ddot \eta &= \frac{\langle\ddot \psi, \psi\rangle_\mathcal{M}}{\eta} + \frac{\|\dot \psi\|_\mathcal{M}^2}{\eta} - \frac{\left(\langle \dot \psi, \psi\rangle_\mathcal{M}\right)^2}{\eta^3}
\\
&\ge \frac{\langle \ddot \psi, \psi\rangle_\mathcal{M}}{\eta}
\\
&=\frac{\langle \nabla U(r_0 + \psi) - \nabla U(r_0),\,\psi\rangle}{\eta}\\
& = \frac{1}{\eta} \int_{0}^{1}\langle\nabla^2 U(r_0+s\psi)\psi,\,\psi\rangle \, \ud s.
\end{split}
\]
For the sake of readability, we define
\[
x_s(t):=r_0(t)+s\,\psi(t) \ \mbox{ and } \ \hat x_s(t)=\frac{x_s(t)}{\|x_s(t)\|_{\mathcal M}}=\frac{r_0(t)+s\,\psi(t)}{\|r_0(t)+s\,\psi(t)\|_{\mathcal M}}.
\]
From \eqref{eq:hessian}, we have
\[
\langle\nabla^2 U(x_s)\psi, \psi\rangle \ge  \frac{\langle\nabla^2\tilde  U(\hat x_s)\psi,\psi\rangle}{\|x_s\|_{\mathcal M}^{2+\alpha}} - \frac{\langle  \nabla\tilde U(\hat x_s), \psi\rangle \langle \hat x_s, \psi\rangle_{\mathcal M}}{\|x_s\|_{\mathcal M}^{2+\alpha}} - \alpha\frac{U(\hat x_s)}{\|x_s\|_{\mathcal M}^{2+\alpha}}\|\psi\|_{\mathcal M}^2.
\]
Since $\hat x_s \to b_m$, and $b_m$ is a local minimum for $\tilde U$, then, for every $\varepsilon>0$, there is $T$ such that for all $t\ge T$
\[
\langle\nabla^2 U(x_s(t))\psi, \psi\rangle \ge -(1-\varepsilon) \frac{\alpha\,U_{min}}{\|r_0\|_{\mathcal M}^{2+\alpha}} \|\psi\|_\mathcal{M}^2 = -(1-\varepsilon)\frac{2\alpha}{(2+\alpha)^2}\frac{\eta^2}{ {t^2}}.
\]

Putting this into the differential inequality for $\eta$, we get
\[
\ddot\eta \ge  -(1-\varepsilon)\frac{2\alpha}{(2+\alpha)^2}\frac{\eta}{t^2}, \ \mbox{ for } \ t \ge T,
\]
that is $\ddot \eta + \mu_\varepsilon \frac{\eta}{t^2} \ge 0$, for $t \ge T$, with
\[
\mu_\varepsilon = (1-\varepsilon)\frac{2\alpha}{(2+\alpha)^2}< \frac{2\alpha}{(2+\alpha)^2}<1/4.
\]
By means of Proposition \ref{prop:max_principle}, with $f=0$, we obtain, for every $\varepsilon>0$,
\begin{equation}
    \label{eq:gamma_eps}
\eta(t) \le C_T\,t^{\sigma(\varepsilon)}, \ t \ge T,
\end{equation}
where $C_T=\frac{\eta(T)}{T^{\sigma(\varepsilon)}}$ and
\[
\sigma(\varepsilon) = \frac{1-\sqrt{1-4\mu_\varepsilon}}{2} = \frac{\alpha+\bar\varepsilon}{2+\alpha},
\]
for $\bar\varepsilon >0$ that tends to $0$ if $\varepsilon \to 0^+$.

\smallskip
\emph{Step 2. Refined estimates on $\eta$.} \ We can go back to the equation for $\eta(t)=\|\psi(t)\|_{\mathcal M}$, as in \textit{Step 1}:
\[
\begin{aligned}
\ddot \eta &\ge \frac{ \int_0^1 \langle \nabla^2 U(x_s)\psi, \, \psi \rangle \, \ud s}{\|\psi\|_\mathcal{M}} \ge \int_0^1 \left(-\alpha\frac{\tilde U(\hat x_s)\|\psi\|_\mathcal{M}^2}{\|x_s\|_\mathcal{M}^{2+\alpha}\|\psi\|_\mathcal{M}} + \frac{\langle\nabla^2 \tilde U(\hat x_s)\psi, \psi\rangle}{\|x_s\|_\mathcal{M}^{2+\alpha}\|\psi\|_\mathcal{M}} + \frac{\langle \nabla\tilde U(\hat x_s), \psi\rangle \langle \hat x_s, \psi\rangle_\mathcal{M}}{\|x_s\|_{\mathcal M}^{2+\alpha}\|\psi\|_\mathcal{M}}\right) \ud s
\\
&=-\alpha\frac{\tilde U(b_m)}{\|r_0\|_\mathcal{M}^{2+\alpha}} \eta + \alpha\|\psi\|_\mathcal{M}\int_0^1 \left(\frac{\tilde U(b_m)}{\|r_0\|_\mathcal{M}^{2+\alpha}} -\frac{\tilde U(\hat x_s)}{\|x_s\|_\mathcal{M}^{2+\alpha}}\right)\,\ud s + \int_0^1 \frac{\langle\nabla^2 \tilde U(\hat x_s)\psi, \psi\rangle}{\|x_s\|_\mathcal{M}^{2+\alpha}\|\psi\|_\mathcal{M}}\,\ud s +\int_0^1 \frac{\langle \nabla\tilde U(\hat x_s), \psi\rangle \langle \hat x_s, \psi\rangle_\mathcal{M}}{\|x_s\|_\mathcal{M}^{2+\alpha}\|\psi\|_\mathcal{M}}\, \ud s
\\
&{\ge -\alpha\frac{U_{min}}{\|r_0\|_\mathcal{M}^{2+\alpha}} \eta + \alpha\|\psi\|_\mathcal{M}\int_0^1 \left(\frac{U_{min}}{\|r_0\|_\mathcal{M}^{2+\alpha}} -\frac{\tilde U(\hat x_s)}{\|r_0+s\,\psi_b\,b\|_\mathcal{M}^{2+\alpha}}\right)\,\ud s + I_2+I_3}
\\
&=-\frac{2\alpha}{(2+\alpha)^2} \frac{\eta}{t^2} + I_1+I_2+I_3,
\end{aligned}
\]
where $\psi_b:=\langle \psi(t),b_m\rangle_{\mathcal M}$.

We claim that $I_1+I_2+I_3=O(t^{-q})$, as $t \rightarrow+\infty$, , for some $q \ge \frac{4+\alpha}{2+\alpha}$.

About $I_1$,  we have
{
\[
\begin{aligned}
 \frac{U_{min}}{\|r_0\|_\mathcal{M}^{2+\alpha}}&-\frac{\tilde U(\hat x_s)}{\|r_0+s\psi_b\,b_m\|_\mathcal{M}^{2+\alpha}}
=
\frac{1}{\|r_0\|_\mathcal{M}^{2+\alpha}}
\left[
U_{min}
-
\frac{\tilde U(\hat x_s)}{
\left( 1 + s^2\frac{\psi_b^2}{\|r_0\|_\mathcal{M}^2} + 2s\psi_b \frac{\langle b_m,r_0\rangle_\mathcal{M}}{\|r_0\|_\mathcal{M}^2} \right)^{\frac{2+\alpha}{2}}}
\right]
\\[0.6em]
& \le \frac{1}{\|r_0\|_\mathcal{M}^{2+\alpha}}
\left[
U_{min}
-
\frac{\tilde U(\hat x_s)}{
\left( 1 + \frac{|\psi_b|^2}{\|r_0\|_\mathcal{M}^2} +  \frac{ 2|\psi_b|}{\|r_0\|_\mathcal{M}} \right)^{\frac{2+\alpha}{2}}}
\right]
\\[0.6em]
&=\frac{1}{\|r_0\|_\mathcal{M}^{2+\alpha}}\left[U_{min} - \tilde U(\hat x_s) \left(1 - \frac{(2+\alpha)}{2}\left(\frac{|\psi_b|^2}{\|r_0\|_\mathcal{M}^2} + \frac{2|\psi_b|}{\|r_0\|_\mathcal{M}}\right)\right)+O\left(\frac{\|\psi\|_\mathcal{M}}{\|r_0\|_\mathcal{M}}\right)\right]
\\
&=\frac{1}{\|r_0\|_\mathcal{M}^{2+\alpha}}\left[U_{min}-\left(U_{min}+s\frac{\langle \nabla^2\tilde U(b_m)(\psi-\psi_b\,b_m),\psi- \psi_b\,b_m\rangle}{2\|r_0\|_\mathcal{M}^2}\right)\left(1-(2+\alpha)\frac{|\psi_b|}{\|r_0\|_{\mathcal M}}\right)\right.
\\
&\quad \quad \quad +\left.O\left(\frac{\|\psi\|_\mathcal{M}}{\|r_0\|_\mathcal{M}}\right)\right]
\\
&=O\left(\frac{|\psi_b|}{\|r_0\|_{\mathcal M}^{3+\alpha}}\right),
\end{aligned}
\]
where we used, in particular, that
\begin{equation}\label{eq:x_hat}
\begin{aligned}
    \hat x_s &= \frac{r_0+s\,\psi}{\|r_0+s\,\psi\|_\mathcal{M}} = \frac{r_0}{\|r_0+s\,\psi\|_\mathcal{M}}+\frac{s\,\psi}{\|r_0+s\,\psi\|_\mathcal{M}}
    \\
    &=b_m+\left(\frac{b_m}{\|b_m+s\,\psi/\|r_0\|_\mathcal{M}\|_\mathcal{M}}-b_m\right) + \frac{s\,\psi}{\|r_0\|_\mathcal{M}} + O\left(\frac{\|\psi\|_\mathcal{M}^2}{\|r_0\|_\mathcal{M}^2}\right)
    \\
    &= b_m + \frac{s}{\|r_0\|_\mathcal{M}}\left(\psi-\psi_b\,b_m\right)+ O\left(\frac{\|\psi\|_\mathcal{M}^2}{\|r_0\|_\mathcal{M}^2}\right).
    \end{aligned}
\end{equation}
}
Therefore, by means of {\em Step 1}, we get
\[
\begin{aligned}
I_1&=\alpha\|\psi\|_\mathcal{M}\int_0^1\left(
\frac{U_{min}}{\|r_0\|_\mathcal{M}^{\alpha+2}}
-
\frac{\tilde U(\hat x_s)}{\|r_0+s\,\psi_b\,b_m\|_\mathcal{M}^{2+\alpha}}\right)\ud s
\\
& =
O\!\left(
\frac{|\psi_b|\|\psi\|_\mathcal{M}}{\|r_0\|_\mathcal{M}^{3+\alpha}}
\right) = O\!\left(
\frac{\|\psi\|_\mathcal{M}^2}{\|r_0\|_\mathcal{M}^{3+\alpha}}
\right) \le C\, t^{\frac{2\alpha}{2+\alpha}+\varepsilon}t^\frac{-2(3+\alpha)}{2+\alpha}= C\,t^{-\frac{6}{2+\alpha}+\varepsilon},
\end{aligned}
\]
for every $\varepsilon>0$.

About $I_2$, we make use of the following fact. For every $V \in C^4$, if $b$ is a local minimum of $V$, then there exists $c>0$ such that, for $\|y\|_\mathcal{M}$ small enough, we have
\[
\langle \nabla^2 V(b+y)\,y,\,y\rangle \ge -c \|y\|_\mathcal{M}^4.
\]

From \eqref{eq:x_hat}, denoting  $y(t)=s\frac{\psi-\psi_b\,b_m}{\|r_0\|_\mathcal{M}}$, we get
\[
\begin{aligned}
\frac{\langle\nabla^2 \tilde U(\hat x_s)\psi, \psi\rangle}{\|x_s\|_\mathcal{M}^{2+\alpha}\|\psi\|_\mathcal{M}} &= \frac{\langle \nabla^2 \tilde U(b_m+y)y,\,y\rangle}{\|x_s\|_\mathcal{M}^{\alpha}\|\psi\|_\mathcal{M}}+s\psi_b^2\frac{\langle \nabla^2 \tilde U(b_m+y) b_m, \, b_m\rangle}{\|x_s\|_\mathcal{M}^{2+\alpha}\|\psi\|_\mathcal{M}}
\\
&\ge -c \frac{\|y\|_\mathcal{M}^4}{\|r_0\|_\mathcal{M}^\alpha \|\psi\|_\mathcal{M}}-C\frac{|\psi_b|^2\,\|y\|_\mathcal{M}}{\|r_0\|_\mathcal{M}^{2+\alpha} \|\psi\|_\mathcal{M}}
\\
&= -c \frac{\|\psi\|_\mathcal{M}^3}{\|r_0\|_\mathcal{M}^{4+\alpha}}-C'\frac{|\psi_b|^2}{\|r_0\|_\mathcal{M}^{3+\alpha}}
\\
&\ge -c \frac{\|\psi\|_\mathcal{M}^3}{\|r_0\|_\mathcal{M}^{4+\alpha}}-C'\frac{\|\psi\|_\mathcal{M}^2}{\|r_0\|_\mathcal{M}^{3+\alpha}}
\\
&\ge -C''\,t^{-\frac{8-\alpha}{2+\alpha}+\varepsilon}-C''\,t^{-\frac{6}{2+\alpha}+\varepsilon}
\\
&\ge -C'''\,t^{{-\frac{8-\alpha}{2+\alpha}}+\varepsilon},
\end{aligned}
\]
for every $\varepsilon>0$. Here, we used {\em Step 1}, and the fact that, for $\|y\|_\mathcal{M}$ small enough,
\[
\langle \nabla^2 \tilde U(b_m+y)b_m,\,b_m\rangle = \int_0^1 \nabla^3 \tilde U(b_m+s\,y)[b_m,b_m,y]\,\ud s\ge -C\|y\|_\mathcal{M}\ge -C' \frac{\|\psi\|_\mathcal{M}}{\|r_0\|_\mathcal{M}}.
\]

About $I_3$, we observe that
\[
\begin{aligned}
    \langle \nabla \tilde U(b_m+y),\,\psi\rangle &= \int_0^1 \langle\nabla^2 \tilde U(b_m+s\,y)y,\,\psi-\psi_b\,b_m \rangle\,\ud s + \int_0^1\langle \nabla^2 \tilde U(b_m+s\,y)y,\,\psi_b\,b_m\rangle\,\ud s 
    \\
    & = \|r_0\|_\mathcal{M}\int_0^1 \langle\nabla^2 \tilde U(b_m+s\,y)y,\,y \rangle\,ds +\int_0^1\int_0^1 \nabla^3 \tilde U(b_m+s\,u\,y)[y,y,b_m]\,\ud s\,\ud u
    \\
    &\ge -c \|r_0\|_\mathcal{M}\|y\|_\mathcal{M}^4 -C \|y\|_\mathcal{M}^2\,|\psi_b|.
\end{aligned}
\]
Hence,
\[
\begin{aligned}
I_3 &\ge  -c \frac{\|r_0\|_\mathcal{M}\|y\|_\mathcal{M}^4}{\|r_0\|_\mathcal{M}^{2+\alpha}\|\psi\|} -C \frac{\|y\|_\mathcal{M}^2\,|\psi_b|}{\|r_0\|_\mathcal{M}^{2+\alpha}\,\|\psi\|_\mathcal{M}}
\\
&\ge -c \frac{\|\psi\|_\mathcal{M}^3}{\|r_0\|_\mathcal{M}^{5+\alpha}}-C'\frac{\|\psi\|_\mathcal{M}^2}{\|r_0\|_\mathcal{M}^{4+\alpha}}
\\
&\ge -C'' t^{-\frac{8-\alpha}{2+\alpha}+\varepsilon}.
\end{aligned}
\]
for every $\varepsilon>0$, by means again of {\em Step 1}.

\smallskip
\emph{Step 3. Conclusion.}
\ Putting together Steps 1--3, we have proved that $\eta$ satisfies
\[
\ddot \eta -\frac{2\alpha}{(2+\alpha)^2}\eta \ge f(t), 
\]
with $f(t)=O(t^{-\frac{8-\alpha}{2+\alpha}+\varepsilon})$, as $t\rightarrow+\infty$, for every $\varepsilon>0$. Hence, an application of Proposition \ref{prop:max_principle} and Remark \ref{rem:corollary} leads to  
\[
\eta(t) \le \frac{\eta(T)}{T^{\theta_-}}t^{\theta_-} + y_p(t)=O(t^{\theta_-})=O(t^\frac{\alpha}{2+\alpha}), \ \mbox{ as } \ t\rightarrow+\infty,
\]
that concludes the proof.

\end{proof}

\end{document}
